% =====================================================================
% Section: Audit of published near-lossless claims.
%
% RESTRUCTURED 2026-08-04 (flagship narrative). This section now runs AFTER
% \S\ref{sec:atlas}, which repairs a live logical defect: the discordance
% imputation below is defined over the atlas, and the atlas used to appear two
% sections later.
% MOVED TO Appendix~\ref{app:audit-method} IN THE SAME COMMIT, NOTHING DELETED:
%   tab:audit-locus, tab:audit-mdd (both whole, every column kept)
% MOVED TO Appendix~\ref{app:prereg:choices}: the K = 5 -> K = 4 R04 history,
%   which that appendix already carried in full as interpretive choice 3.
%
% PRIMARY SOURCES for every number in this section:
%   results/audit_verdicts_rev3.csv     (AUTHORITATIVE, per-claim, machine-readable;
%                                        job 11591245, sha c85d6f8a...b150082b,
%                                        single run, chmod 0444 -- never re-run)
%   docs/audit_claim_table.csv          (frozen inputs, sha256 in the doc)
%   docs/AUDIT_REGISTRATION_2026-07-15.md (frozen protocol, §§3-5, + Amendment 2)
%   docs/AUDIT_SOURCE_VERIFICATION_2026-07-31.md (full-text review of all 17)
%
% REV-3 SUPERSESSION, 2026-07-31. Amendment 2 (signed, commit ab279b2) is in
% force. It (a) makes the registered UNIFORM 2 pp margin the primary yardstick
% for every claim, (b) forbids describing any quantity derived from a source's
% own reported results as that source's own/stated/declared margin, and (c)
% reopens §§3.1-3.2 for eligibility only, excluding R10. Consequences for this
% section, all carried below:
%   K = 4 of 12 "underpowered for their own assertion"  ->  1 of 11 below the
%     approximate planning threshold at 2 pp, and that one is imputation-
%     sensitive, so 0 of 11 are below it throughout the atlas IQR.
%   denominators 17 -> 16 eligible.
%   the 2.0x-12.9x shortfall range is WITHDRAWN, not recomputed: ten of eleven
%     claims have no shortfall at all at 2 pp, so there is no range to state.
% docs/AUDIT_VERDICTS_2026-07-20.md is the SUPERSEDED record (rev-1 atlas AND
% pre-Amendment-2 margins). Do not "restore" any value from it.
%
% SUPERSESSION NOTICE -- do not "fix" this section against the older doc.
% docs/RESULTS_2026-07-15_ATLAS_AUDIT.md line 54 states "only 1 of 17 claims
% has per-item outputs released". That line is SUPERSEDED by the 2026-07-20
% verdicts document, which records 0 yes / 3 partial / 14 no. This paper uses
% the rev-3 denominator; the three "partial" sources are the ones the older line
% was gesturing at.
%
% DISCREPANCY -- CLOSED 2026-07-26. THE FROZEN TABLE GOVERNS.
% docs/RESULTS_2026-07-15_ATLAS_AUDIT.md line 48 describes the frozen table as
% "7 method papers, 8 official model cards/blog, 2 vendor docs" (7/8/2). The
% frozen table itself gives F1=7, F2=7, F3=3, re-verified 2026-07-26 by counting
% the `frame` column of docs/audit_claim_table.csv (7/7/3, 17 rows total).
% docs/audit_claim_table.csv is a FROZEN artifact and is authoritative over the
% older narrative doc, so the prose below (7/7/3) is correct as written.
%
% REV-2 POINTER CORRECTION, 2026-07-27, RETAINED AS HISTORY. Every SOURCE
% comment in this tree that cited results/audit_verdicts.csv was repointed to
% the rev-2 file; the typeset values were already rev-2, so only the provenance
% pointers had been wrong, and they were instructing future sessions to
% "correct" correct values back to superseded rev-1 ones. That hazard is why the
% rev-2 -> rev-3 repointing was done comment by comment rather than by a blind
% replacement. Rev-2 -> rev-3 is a different kind of change: the yardstick
% itself moved (Amendment 2), so verdicts DID move.
% =====================================================================

\section{What published equivalence claims actually report}

% SOURCE: results/audit_verdicts_rev3.csv (job 11591245), computed once under
% Amendment 2 over the frozen claim table (sha 842b9756...) and the rev-2 atlas
% (sha b7cffc52...), both DECLARED on the command line, neither defaulted.
\label{sec:audit}

\subsection{What was audited, and what was not}

% STAGE A of the audit compression, 2026-08-06. Protected qualifications 1, 2,
% 3, 12, 13 and 18 all live in this subsection and are carried below.
% RELOCATED, NOT DELETED: the full-text review behind the R10 ruling and the
% source provenance detail (content hashes for 16 of 17, R13 empty) are in
% Appendix~\ref{app:audit:eligibility}.
% NEVER describe the two extraction passes as inter-rater or independent human
% verification. NEVER collapse 17 and 16 into one number.
The audit protocol was frozen on 2026-07-15, before any per-claim power
computation was run, and the claim list itself was frozen separately before any
verdict was computed.\footnote{%
% SOURCE: docs/AUDIT_REGISTRATION_2026-07-15.md (status line, §3.4);
% freeze provenance reconfirmed on the rev-3 command line (job 11591245).
Registration frozen at commit \texttt{b74fd58}; claim table frozen at commit
\texttt{715a7ce} with its sha256 recorded in the verdicts document.}
Sources were enumerated exhaustively within three fixed frames: method papers
(F1), official quantized model cards (F2), and inference-stack vendor blogs and
documentation (F3). A source enters the pool if it asserts, in prose or a table
caption, that a compressed model's benchmark quality is
equivalent-or-negligibly-different from its uncompressed baseline, under a
trigger vocabulary fixed in advance. Every claim meeting the criterion is
audited; there is no discretionary sub-selection.

% SOURCE: docs/audit_claim_table.csv column `frame`: F1=7, F2=7, F3=3.
% BOTH DENOMINATORS ARE LOAD-BEARING AND MUST NOT BE COLLAPSED.
The frozen table contains \AuditFrozenCandidates{} candidate claims: 7 method
papers, 7 official model cards or blogs, and 3 vendor documents. A full-text
review of every source, conducted after the first verdicts were computed, found
that one candidate's recorded quotation appears nowhere in its source, having
been composed from a table cell. The registered rule requires the assertion to
appear in prose or a table caption, so \AuditIneligibleClaim{} is excluded by
applying that rule as registered, leaving \textbf{\AuditEligible{}} eligible
sources; every count in this section is over those \AuditEligible{}, and the
exclusion lowers no count reported here
(Appendix~\ref{app:audit:eligibility}).

% ADDED 2026-07-31, PROTECTED. The body must say WHAT the passes were.
Each claim's fields were extracted by two blinded automated passes, separate
language-model agent sessions with the second given no access to the first
pass's output, followed by human reconciliation of every disagreement before the
verdict stage (\S\ref{sec:prereg}). Separate sessions prevent the second pass
from copying the first; they do \emph{not} make the two passes statistically
independent, because both share a common model prior, and nothing here is
inter-rater agreement. Appendix~\ref{app:extraction} specifies the procedure and
reports what can and cannot be measured after the fact.

% SOURCE: docs/AUDIT_REGISTRATION_2026-07-15.md §4 closing paragraph.
% PROTECTED: evidential sufficiency, not truth. No claim is called false. The
% prevalence disclaimer is stated in full at the section close as well.
\textbf{No audited claim is described as false, and none of our findings implies
that any audited model is in fact degraded.} The audited property is the
\emph{evidential sufficiency of the reported evaluation}: whether the evaluation
offered in support of an equivalence claim was large enough to have detected the
difference the claim pronounces negligible. A claim can be perfectly correct and
still be unsupported by the evaluation offered for it, and several of the claims
below are very probably correct. Nor is any proportion reported here a
prevalence estimate for the field. What is missing from the audited sources is a
reporting standard: this section measures the gap such a standard would close,
and \S\ref{sec:certification} supplies the instrument.


\subsection{What the sources declare}
\label{sec:audit:taxonomy}

% ADDED 2026-07-31, carrying AUDIT_REGISTRATION Amendment 2 (signed, commit
% ab279b2) and docs/AUDIT_SOURCE_VERIFICATION_2026-07-31.md.
% SOURCE: results/audit_verdicts_rev3.csv columns eligible, eligibility_basis,
% margin_category, evidence_form. Job 11591245.
%
% DO NOT report the marginals of margin_category alone. "0 formal / 12 informal
% / 4 unquantified" is arithmetically right and rhetorically wrong: 7 of the 12
% category-2 claims contain NO NUMBER AT ALL, so calling them informal *margins*
% reads as contradicting the 0-of-16 headline in the same paragraph. The
% cross-tab is the honest presentation.
%
% STAGE A, 2026-08-06: the R10 eligibility ruling moved UP into
% \S\ref{sec:audit} (design and scope), where both denominators are now
% introduced together, and the full-text review behind it moved to
% Appendix~\ref{app:audit:eligibility}. Nothing was deleted.
Across the \AuditEligible{} eligible sources we distinguish two independent
properties: what form the source's evidence takes, and whether the reported
information suffices to assess it numerically at all.

\begin{table}[!t]
\centering
\caption{What the 16 eligible sources declare. Columns are the form of the
evidence offered; rows are whether the source reports enough for a numerical
assessment under the registered paired-outcome framework. \emph{Prospective
numerical decision margin} means a tolerance stated independently of the
observed result: a threshold that could have been written down before the
evaluation ran. That column is empty.}
\label{tab:audit-taxonomy}
% SOURCE: results/audit_verdicts_rev3.csv, cross-tabulating margin_category
% (2=assessable-informal, 3=unquantified) against evidence_form. Verified to
% sum to 16 on both margins: 7+5=12, 3+1=4, 7+3=10, 5+1=6.
% WIDTH FIX 2026-08-05: was {lccc} with three long free-running header phrases,
% which ran 133pt past the measure. Fixed-width centred header columns wrap the
% phrases instead. No cell value changed, and the empty prospective column that
% is this table's whole evidential point is untouched.
\setlength{\tabcolsep}{4pt}
\begin{tabular}{l>{\centering\arraybackslash}p{2.2cm}>{\centering\arraybackslash}p{2.4cm}>{\centering\arraybackslash}p{2.2cm}}
\toprule
 & Prospective numerical decision margin & Retrospective numerical description of a result & Qualitative language only, no number \\
\midrule
Numerically assessable     & 0 & 5 & 7 \\
Not numerically assessable & 0 & 1 & 3 \\
\midrule
\textbf{Total}             & \textbf{0} & \textbf{6} & \textbf{10} \\
\bottomrule
\end{tabular}
\end{table}

% SOURCE: docs/AUDIT_SOURCE_VERIFICATION_2026-07-31.md §§6-7.
% STAGE B, 2026-08-06: the sweep's contents (parity and percentage point absent
% from the corpus, every "margin" a layout artefact or the "by a large margin"
% idiom, the one vendor explicitly declining to fix a threshold) moved to
% Appendix~\ref{app:audit:sweep}. Nothing was deleted.
% PROTECTED: the zero, the 10, and that the remaining 6 cite a MEASURED OUTCOME
% rather than a declared tolerance. Do not report the margin_category marginals
% alone: 7 of the 12 category-2 claims contain no number at all, so calling them
% informal *margins* reads as contradicting the zero in the same paragraph.
\textbf{No audited source declares a prospective numerical equivalence margin}
(Table~\ref{tab:audit-taxonomy}). \AuditXtabQualTotal{} of the
\AuditEligible{} make the claim in purely qualitative terms (``negligible'',
``matches'', ``practically no accuracy decrease''), with no number attached to
the assertion at all. The remaining \AuditXtabRetroTotal{} cite a number, and in
every case it is a \emph{measured outcome}: a recovery percentage or an achieved
benchmark score, computed after the evaluation rather than set before it. This
is established over complete source text, not the quoted sentence: every source
was archived and searched in full, including tables, captions, footnotes and
appendices, for a registered vocabulary of tolerance language
(Appendix~\ref{app:audit:sweep}).

The consequence for this section is structural. An equivalence claim with no
declared margin cannot be evaluated against its own standard, because it has
set none; it can only be evaluated against a standard supplied from outside.
Everything below therefore uses the margin registered in advance
(\S\ref{sec:audit:rules}), and no quantity derived from a source's own reported
results is described anywhere in this paper as that source's margin.

\subsection{Where the claim is written}
\label{sec:audit:locus}

% STAGE B, 2026-08-06. The six-card comparison, the tier description and the
% retention judgement for the two boundary cases moved WHOLE to
% Appendix~\ref{app:audit:locus}, which already held tab:audit-locus.
% RETAINED HERE, PROTECTED: the consequence for the denominator (qualification
% 13's floor-not-census reading) and the author-re-verification disclosure
% (qualification 18). This locus review was a second automated pass by the same
% class of tool that produced the record it checked. It is NOT independent
% verification and must never be described as dual coding or inter-rater
% agreement anywhere in this paper.
% STILL OPEN (defect D6): the 3/2/1 tier counts and the six-card denominator are
% hand-typed in the appendix. They are not emitted by gen_denominator_macros.py,
% which is another session's active file. Do not add a second generator.
A second question turns out to matter as much as what form an assertion takes:
\emph{where in the document it is written}, and whether it is written in words at
all. Across six model cards from one publisher with the underlying evidence held
constant, the claim appears in prose in three, as a bare juxtaposition of two
scores in two, and only as a table cell in one
(Table~\ref{tab:audit-locus}, Appendix~\ref{app:audit:locus}). The registered
inclusion rule is keyed to a property that varies independently of the evidence,
so the claims do not become weaker as the locus moves from sentence to cell:
they stop being written down. The consequence is that
\AuditFrozenCandidates{} is a floor on the population and not a census of it,
with the shortfall concentrated in exactly the sources that state their claims
least explicitly. This bounds what the counts here can mean without weakening
them, since every reported quantity is computed over the sources that were
captured, and it means the true rate of unbounded equivalence claims in this
literature is understated here rather than overstated.

The locus review is author re-verification against the archived sources, not
independent verification by a second coder; the project has no second human, and
Appendix~\ref{app:extraction} states what that does and does not license. Every
archived source, its content hash, and the retrieval script needed to reproduce
these locations are released with the paper, so the classification is checkable
by any reader without re-fetching anything.

\subsection{Availability of per-item outputs}
\label{sec:audit:v3}

The most actionable finding requires no statistics at all.

% SOURCE: results/audit_verdicts_rev3.csv column v3_per_item_outputs over the 16
% ELIGIBLE rows: 0 "yes", 3 "partial" (R08, R15, R16), 13 "no". R10's row is
% "no" and is excluded here, which is why the count is 13 and not 14.
%
% THE QUALIFIER IS LOAD-BEARING AND MUST STAY IN THE SAME SENTENCE. Three
% sources DO release per-item outputs -- for other task suites. "None released
% per-item outputs" is false without "task-matched".
\begin{quote}
\textbf{\AuditPerItemTaskMatched{} of the \AuditEligible{} eligible sources
release \emph{task-matched} per-item outputs}, meaning the item-level results,
for the tasks the equivalence claim is actually about, that a third party would
need to rerun the paired comparison. The tally is \AuditPerItemTaskMatched{}
\emph{yes}, \AuditPerItemOtherTaskOnly{} \emph{partial}
(\AuditPerItemOtherTaskClaims{}), \AuditPerItemNone{} \emph{no}.
\end{quote}

% STAGE B, 2026-08-06: the suite-by-suite detail for the three partial sources
% and the repair-cost argument moved to Appendix~\ref{app:audit:v3detail}.
% PROTECTED AND RETAINED: the qualifier that the three DO release per-item
% outputs, for other suites. "None released per-item outputs" is false without
% "task-matched", and the zero above is meaningless without this sentence.
\AuditPerItemOtherTaskClaims{} come closest: they release per-item outputs for
other suites, but not for the tasks their audited equivalence claim is about, so
the released artifacts and the asserted claim do not intersect
(Appendix~\ref{app:audit:v3detail}). This finding is more actionable than any
power calculation because the two failures have different repair costs.
Underpowering is fixable by evaluating more items, and \S\ref{sec:certification}
says how many. Irreproducibility is not fixable downstream at all: with no
per-item outputs, nobody outside the releasing organisation can run the paired
test at \emph{any} sample size, cannot compute churn, and cannot check the
arithmetic. Per-item outputs for a benchmark of a few thousand items are a file
of a few megabytes, and the gap between that cost and its consequence is the
single clearest reporting-standards recommendation this paper makes.

\subsection{Verdict rules}
\label{sec:audit:rules}

% STAGE C, 2026-08-06. The three estimators, eq:tost-n and the one-sided/90%
% convention moved WHOLE to Appendix~\ref{app:audit:verdictrules}; the
% tier-by-tier imputation detail was already in
% Appendix~\ref{app:audit:imputation}. eq:tost-n has no reference outside this
% file, so it moved with the equation.
% RETAIN VERBATIM: the margin sentence below is the logical hinge of the section
% (mandatory qualification 5) and has now survived two compressions unaltered.
Because no source declares a margin, \textbf{the applicable margin throughout is
the uniform \AuditMarginPP{}\,pp registered in advance}; a source's reported
deltas are outcomes of its evaluation, not a decision rule it adopted, and are
never used as one.

For each claim, at its reported or rule-imputed sample size, the frozen protocol
computes three quantities: \emph{V1}, the minimum detectable difference under
the paired-flip model; \emph{V2}, the number of items TOST would require at the
applicable margin; and \emph{V3}, whether per-item outputs were released. A
claim is recorded as \emph{below the approximate planning threshold} iff its
reported $n$ falls below the V2 requirement.
Appendix~\ref{app:audit:verdictrules} gives the estimators and the one-sided
$z$ convention, under which the requirement is that a 90\% two-sided interval
fall inside the margin.
% WORDING RULE (Amendment 2 + carry checklist §8): a PROSPECTIVE PLANNING
% quantity, not a retrospective diagnosis. Do not reintroduce a definitive
% underpowered verdict here.
That requirement answers a question asked \emph{before} an evaluation is run,
namely how many items it would need, so falling below it is a statement about
how the evaluation was sized, not a diagnosis applied to its result.

% RELOCATED 2026-07-27: tier-by-tier match counts and the None-descent rule are
% in Appendix~\ref{app:audit:imputation}. Nothing was deleted.
The discordance rate $p_d$ is reported by no source, so it is imputed from the
atlas (\S\ref{sec:atlas}) by matching the nearest (method family, bit width,
benchmark) cell, most-specific tier first, taking the \emph{median} over the
first non-empty tier because per-cell discordance is right-skewed
(Appendix~\ref{app:audit:imputation}). Because $p_d$ is imputed and never
reported, each verdict is additionally recomputed at the first and third
quartiles of the same atlas cells that supplied its median.
% CHRONOLOGY DISCLOSURE, REQUIRED (carry checklist §8). Added 2026-07-31, AFTER
% the point-imputation result was seen. Saying so is not optional and it must
% not be phrased so as to imply it was planned.
This quartile sweep was added after the point-imputation results had been seen,
and is reported as post-hoc sensitivity, not as a registered analysis.
% Monotonicity, not cell-counting, is what makes the interval claim valid.
Because $n_{\mathrm{req}}$ is increasing in $p_d$, the two quartiles bracket the
whole interval: a claim above threshold at $Q_3$ is above it throughout, and one
below threshold at $Q_1$ is below it throughout. The interval statement is
established at its endpoints, not inferred from the scatter of observed cells.

\subsection{Results}
\label{sec:audit:results}

% SOURCE: results/audit_verdicts_rev3.csv columns eligible, indeterminate,
% v2_underpowered_paired_2pp, robustness. Counts over the 16 ELIGIBLE rows.
% DENOMINATOR LEDGER (carry checklist §9): 17 - 1 = 16; 11 + 5 = 16;
% 0 + 1 + 10 = 11. Every one of these is checkable in the CSV.
Of the \AuditEligible{} eligible claims, \AuditNotAssessable{} are \textbf{not
numerically assessable}: their reporting does not support a verdict under the
registered framework, leaving \AuditAssessable{} assessable claims. Of those
\AuditAssessable{}, \AuditBelowThresholdAtMedian{} falls below the
approximate planning threshold at the registered \AuditMarginPP{}\,pp margin
under the median discordance imputation, and that one claim's classification
does not survive the sensitivity analysis.

% SUPERSEDED TABLE, DO NOT RESTORE. This replaces tab:audit-underpowered, whose
% four rows were flagged at their own reported deltas and are ALL robustly above
% the threshold at 2 pp. The label itself was a stale claim, so it was renamed
% rather than recaptioned.
% A one-row successor table naming only R01 was considered and REJECTED: a table
% built around the single flag looks constructed to produce it. The main text
% carries the three-way classification; all 11 rows are in the appendix.
\begin{table}[!t]
\centering
\caption{Sensitivity classification of the 11 numerically assessable claims at
the registered 2\,pp margin. Each claim's required sample size is recomputed at
the first and third quartiles of the atlas cells that supplied its median
discordance; because required $n$ increases monotonically in the discordance
rate, the quartiles bracket the whole interval. Per-claim values are in
Appendix~\ref{app:audit-table}.}
\label{tab:audit-sensitivity}
% SOURCE: results/audit_verdicts_rev3.csv column `robustness` over the 11
% eligible determinate rows, via the generated \Audit* macros in
% paper/audit_denominators.tex: "robustly above threshold" =
% \AuditAboveThroughout, "imputation-sensitive" = \AuditChangesWithinIQR (R01),
% "robustly below threshold" = \AuditBelowThroughout.
% DEFECT D5, FIXED 2026-08-04: these three cells were HAND-TYPED as 10, 1, 0.
% No section of this paper types an audit count; that rule exists because the
% retired "4 of 12" headline outlived the verdicts that produced it.
\begin{tabular}{lr}
\toprule
Classification across the atlas-IQR interval & Claims \\
\midrule
Above the planning threshold throughout          & \AuditAboveThroughout{} \\
Changes classification within the interval       & \AuditChangesWithinIQR{} \\
Below the planning threshold throughout          & \AuditBelowThroughout{} \\
\midrule
\textbf{Total assessable}                        & \textbf{\AuditAssessable{}} \\
\bottomrule
\end{tabular}
\end{table}

% STAGE C, 2026-08-06: the reversal-point and quartile arithmetic moved to
% Appendix~\ref{app:audit:r01}; the detection-direction reading of
% tab:audit-mdd moved to Appendix~\ref{app:audit:fullrobustness}; the reasoning
% behind the two sensitivity directions was already in
% Appendix~\ref{app:audit:robustness}.
% THE 43.6% IS DESCRIPTIVE. It is the share of reference cells lying below the
% reversal point. It is NOT a posterior probability, NOT a confidence level and
% NOT a p-value, and the atlas cells are correlated (shared model pairs,
% benchmark families, sources and infrastructure), so they are not independent
% draws from anything. Never attach an inferential reading to it. The number and
% its framing stay in ONE sentence.
The single flagged claim is \AuditSensitiveClaim{}, the GPTQ paper, which
reports $n = \AuditSensitiveN{}$ against a requirement of
$\AuditSensitiveNReq{}$ items at an imputed discordance rate of $0.130$. Its
classification reverses at a rate of $0.1189$, and
$\AuditSensitiveCellsBelow{}$ of the $\AuditSensitiveCellsTotal{}$ atlas cells
supplying the imputation lie below that point
(Appendix~\ref{app:audit:r01}). This is a sensitivity-dependent planning flag,
not a verdict: the $\AuditSensitiveCellsPct{}\%$ is a descriptive share of
reference cells, not a probability that the claim is underpowered, and those
cells share models, benchmarks and infrastructure and are not independent
observations.

No assessable claim falls below the threshold throughout the interval.
Table~\ref{tab:audit-sensitivity} is therefore best read as a statement about
resolution, not about error: at the margin a reporting standard would plausibly
impose, ten of these eleven evaluations were sized to carry the claim made on
them, and the eleventh cannot be classified without committing to a discordance
rate its authors never reported. Nothing here says any model is degraded.

% CARE -- THE TWO SENSITIVITY DIRECTIONS ARE DIFFERENT FACTS AND MUST NEVER BE
% MERGED INTO ONE SENTENCE. Margin sensitivity is instability with respect to
% the yardstick; imputation sensitivity is instability with respect to an input
% the sources did not report. They coincide on R01, which is a coincidence.
The same conclusion appears in the detection direction, where the
independent-binomial columns are uniformly worse by roughly a factor of two
(Table~\ref{tab:audit-mdd}, Appendix~\ref{app:audit:fullrobustness}); an audit at
a \AuditMarginPP{}\,pp margin is a weak test and should be understood as one.
Second, \textbf{R04 and R14 carry no verdict}; their numbers are computable and
are retained in the released CSV for transparency, but they are excluded from
the assessable set (\S\ref{sec:audit:indeterminate}), and the appendix table
italicises them for that reason. Finally, \AuditSensitiveClaim{} is unstable in
two independent directions that must not be read as one finding: it is
\emph{margin-sensitive}, flipping across the registered
1\,pp\,$\to$\,3\,pp sweep, so it turns on where the yardstick is set, and
separately \emph{imputation-sensitive}, flipping across the interquartile range
of the atlas cells supplying its discordance rate, so it also turns on a
quantity its source never reported. No other assessable claim is sensitive in
either direction (Appendix~\ref{app:audit:robustness}).

\subsection{Claims that cannot be assessed}
\label{sec:audit:indeterminate}

% SOURCE: results/audit_verdicts_rev3.csv columns indeterminate,
% indeterminate_kind, indeterminate_reason, over the 16 eligible rows.
% DENOMINATOR: 5 of 16, not of 17. R10 is eligible-excluded, not indeterminate.
% THE 4 + 1 SPLIT IS MANDATORY AND MUST NOT COLLAPSE INTO ONE NUMBER.
% STAGE C, 2026-08-06: the per-claim blockers for R02, R11, R13 and R14 and the
% retained-components list moved to Appendix~\ref{app:audit:indeterminate}. The
% R14 trap paragraph below STAYS IN THE BODY (mandatory qualification 12) and is
% the body's only statement of it.
% NEVER write "incompatible with a paired framework" for R04: CIDEr supports
% paired resampling perfectly well; it is the flip model that does not apply.
\AuditNotAssessable{} of the \AuditEligible{} eligible claims cannot be assessed
numerically, in two kinds, each recorded with exactly one primary blocker.
\AuditNotAssessableInsufficient{} are cases of \emph{insufficient reporting},
where a registered input is genuinely absent: no sample size, no baseline, or
headline equivalence evidence existing only as a chart image with no extractable
numbers. The remaining \AuditNotAssessableOutsideFramework{},
\AuditOutsideFrameworkClaim{}, is \emph{outside the registered calculation}: it
reports enough, but about a quantity our registered \emph{binary paired-outcome}
model cannot score (\S\ref{sec:audit:r04}). The per-claim blockers are in
Appendix~\ref{app:audit:indeterminate}. Every non-assessable claim retains
whatever components its available inputs support, reported as supplementary
transparency only and \textbf{never verdict-bearing}.

% R14 IS THE TRAP THIS PARAGRAPH EXISTS TO CLOSE (carry checklist §10). Its
% imputed n = 728 sits just under the 742 that the 2 pp calculation would
% require, so a reader who sees the numbers and not the blocker will assume it
% was quietly dropped for being awkward. Say the number, then say why it does
% not count. A regression test enforces the same thing in the toolkit: an
% assessable = false row can never carry a threshold verdict.
% THIS PARAGRAPH STAYS IN THE BODY (mandatory qualification 12). The italicised
% no-verdict rows of tab:audit-mdd moved to the appendix in the same commit, so
% this is now the body's only statement of it.
R14 deserves a word, because its numbers look assessable and are not. Its
imputed sample size of $728$ falls just short of the $742$ the 2\,pp calculation
would require, close enough that a reader might suspect it was set aside for
being inconvenient. It is set aside because the missing baseline is a registered
input: without it there is no paired comparison to size, and the $742$ is a
requirement for a calculation that cannot be performed. The proximity is a
coincidence with no bearing on the count, and R14 does not enter the threshold
tally in either direction.

\emph{Outside the registered calculation} (\AuditNotAssessableOutsideFramework{}).
\AuditOutsideFrameworkClaim{} (AWQ) reports enough, but about a quantity our
registered \emph{binary paired-outcome} model cannot score; see below. Every
non-assessable claim retains whatever components its available inputs support,
listed in the CSV column \texttt{determinate\_components} and reported as
supplementary transparency only, never verdict-bearing: R13 retains V2 (the
paired standard deviation depends on discordance, not on baseline accuracy), R14
retains V1 (paired) and V2, and R04 retains V1 and V2 computed on a substituted
benchmark.

This category is itself a result. Two of the four insufficiently-reported claims
are among the most-cited results in the field, and their headline equivalence
evidence is a chart image with no extractable numbers. That a mechanical,
pre-registered audit cannot evaluate them is precisely the reporting-standards
problem this paper exists to address: the claim may well be true, but it has been
placed beyond the reach of checking.

\subsection{R04: outside the registered calculation}
\label{sec:audit:r04}

% SOURCE: results/audit_verdicts_rev3.csv row R04 columns indeterminate_kind,
% indeterminate_reason, notes.
% WORDING RULE (carry checklist §15): R04 is outside OUR REGISTERED BINARY
% PAIRED-OUTCOME CALCULATION. Do NOT write "incompatible with a paired
% framework" -- CIDEr supports paired resampling perfectly well; it is the
% flip model, not pairing, that does not apply.
% K-HISTORY RELOCATED 2026-08-04 (operation 14) to
% Appendix~\ref{app:prereg:choices}, interpretive choice 3, which already
% carried the K = 5 -> K = 4 transition and the 38.3x figure in full. Nothing
% was deleted; the body keeps the ruling and its cost.
R04 is recorded because excluding it removed what was, at the time, the audit's
largest single number. The AWQ paper's qualifying sentence, the one that meets
the frozen §3 inclusion trigger, asserts negligible loss on \emph{COCO
CIDEr}, a generation metric that assigns a graded score to a caption and has no
per-item correct/incorrect state. V1 and V2 are flip-model quantities defined on
$d_i \in \{-1,0,+1\}$, and there is no discordance rate to impute because there
is no per-item accuracy state to be discordant about. This is a limitation of
\emph{our registered calculation}, not of paired analysis: a graded metric
supports paired resampling on the same items perfectly well, and a certificate
for it would be a straightforward extension. It is simply not the model we
registered.

% STAGE D, 2026-08-06: the overruled first-pass GSM8K computation moved to
% Appendix~\ref{app:audit:r04detail}. The body keeps the ruling, that the number
% is not claimed as a result, and its cost.
Excluding it was not free. The first extraction pass had scored R04 on GSM8K,
the source's own accuracy benchmark, and reported it as the audit's largest
shortfall; that reading was overruled because computing a TOST requirement on
GSM8K audits a sentence the source wrote about a different benchmark, in
different and non-trigger language. The computation is retained in the released
CSV as a labelled transparency column, and \textbf{it is not claimed anywhere in
this paper as an audit result} (Appendix~\ref{app:audit:r04detail}). An audit's
currency is unimpeachability, and it is not spent on its own largest number.

% QUALIFICATION 17, PLACED IN THE BODY 2026-08-06. The K sequence is the
% section's canonical self-correction and MUST NOT live only in an appendix. It
% was previously in this file's header comment and in app:prereg:choices, so no
% reader of the paper could see it; the compression ledger caught that. The
% arithmetic of each step stays in the appendix; the sequence itself is here.
% SOURCE: sections/appendix_prereg_detail.tex, interpretive choice 1.
The count of claims below the planning threshold moved four times before it
settled, and the order in which the numbers arrived is itself part of the
disclosure. A first pass applying the registered \AuditMarginPP{}\,pp margin
uniformly returned $K = 1$ of 12. Re-reading the frozen label produced $K = 5$;
the R04 ruling above moved one claim out of the determinate set, giving the
$K = 4$ that was reported for ten days. Amendment~2 then returned the analysis
to the uniform registered margin, where the rev-3 recomputation gives
\textbf{\AuditBelowThresholdAtMedian{} of \AuditAssessable{}} assessable claims
below the threshold, the same claim the very first pass had flagged. The
claim-derived margins and the shortfall range computed under the superseded
reading are \textbf{withdrawn, not recomputed}, and are non-verdict-bearing
wherever they still appear as history
(Appendices~\ref{app:audit:history} and~\ref{app:prereg:choices}).

% TRIMMED 2026-07-27 (phase 4, low-cost). The "Establishes" half restated
% findings (i)-(iv) from \S\ref{sec:audit:v3} and \S\ref{sec:audit:results}
% verbatim; deleted as redundancy. The "Does not establish" half is what a
% reader must not overread and is kept whole, as is the granularity caveat.
% MERGED 2026-08-04: this was a separate closing subsection and is now the
% section's close (prose rule 6). Every clause is retained, including the only
% statement in the body of what "robust" means.
The findings above hold at the reported sample sizes and under the frozen
protocol as amended, and they are claims about evidence, not about models. This section \textbf{does not establish}: that any audited model is
degraded; that any audited claim is false; that the claims' authors reached a
wrong conclusion; or that any claim is underpowered in a sense that survives the
sensitivity analysis, and none is.
% SCOPE RULE (carry checklist §8 and §15). Two overreadings to block explicitly:
% (a) generalising 16 audited sources to a prevalence claim about the field,
% (b) reading "robust" as anything wider than the atlas-IQR interval actually
% swept. Both are cheap to write and expensive to retract.
% THE SECOND OF THOSE IS A SINGLE POINT OF FAILURE: mandatory qualification 9
% has no other home in the body. Do not delete or weaken the sentence beginning
% "Where a classification is called robust".
Nor does it establish a prevalence: \AuditEligible{} eligible sources within
three frozen frames are not a sample from which the field's reporting practice
can be estimated, and no proportion reported here should be read as one. Every
numerical result here also rests on a discordance rate that no source reported
and that is imputed from the atlas, which is why each verdict is recomputed
across the interquartile range of the cells supplying it and why the one flagged
claim is reported as sensitivity-dependent rather than as a finding. Where a
classification is called robust, that means only that it holds throughout the
interquartile range of the atlas cells supplying its discordance rate, not
across every plausible model of discordance. It also does not
establish anything about sources outside the three frozen frames, and it
operates at claim level and not at claim\,$\times$\,benchmark granularity,
because that is the granularity of the frozen claim table
(Appendix~\ref{app:prereg:choices}, interpretive choice~4).
